\par\medskip
\hypertarget{native-075-article057}{}\pdfbookmark[2]{57.}{article-057}
\begin{center}57.\end{center}
Der Abstand eines Himmelskörpers von einer beliebigen durch die Sonne gehenden Ebene lässt sich also auf die Form $kr\sin(v + K)$ bringen. Dabei bezeichnet $v$ die wahre Anomalie, $k$ ist der Sinus der Neigung der Bahn gegen diese Ebene und $K$ der Abstand des Perihels vom aufsteigenden Knoten der Bahn in dieser Ebene. Soweit man die Lage der Bahnebene, die Lage der Apsidenlinie in ihr und die Lage der Bezugsebene für die Abstände als konstant ansehen kann, sind auch $k$ und $K$ konstant. Häufiger wird man dieses Verfahren jedoch in Fällen anwenden, in denen zumindest die dritte Annahme unzulässig ist, selbst wenn man die Störungen vernachlässigt, die die ersten beiden Annahmen stets in gewissem Maße beeinträchtigen. Dies ist immer dann der Fall, wenn man die Abstände auf den
\SourceLeaf{076}Äquator oder auf eine Ebene bezieht, die den Äquator bei gegebener Rektaszension rechtwinklig schneidet. Die Lage des Äquators ändert sich infolge der Präzession der Äquinoktien und, sofern seine wahre Lage anstelle der mittleren gemeint ist, auch infolge der Nutation. In diesem Fall unterliegen daher auch $k$ und $K$ Änderungen, die freilich langsam sind.
Diese Änderungen lassen sich mit leicht herzuleitenden Differenzialformeln berechnen. Der Kürze halber genügt es hier, die differenziellen Änderungen von $i'$, $\AscendingNode'$,
$\Delta$ anzugeben, soweit sie von den Änderungen von $\AscendingNode - n$
und $\epsilon$ abhängen.
\[
\mathrm{d}i' = \sin\epsilon \sin\AscendingNode'\, \mathrm{d}(\AscendingNode - n) - \cos\AscendingNode'\, \mathrm{d}\epsilon
\]
\[
\mathrm{d}\AscendingNode' = \frac{\sin i \cos\Delta}{\sin i'}\, \mathrm{d}(\AscendingNode - n) + \frac{\sin\AscendingNode'}{\operatorname{tang} i'}\, \mathrm{d}\epsilon
\]
\[
\mathrm{d}\Delta = \frac{\sin\epsilon \cos\AscendingNode'}{\sin i'}\, \mathrm{d}(\AscendingNode - n) + \frac{\sin\AscendingNode'}{\sin i'}\, \mathrm{d}\epsilon.
\]
Geht es im Übrigen nur darum, mehrere Positionen eines Himmelskörpers in Bezug auf solche veränderlichen Ebenen innerhalb eines mäßig langen Zeitraums, etwa eines Jahres, zu berechnen, so ist es meist am bequemsten, die Größen $a$, $A$, $b$, $B$, $c$, $C$ für zwei Epochen direkt zu berechnen, zwischen denen diese Positionen liegen. Aus diesen Werten lassen sich dann ihre Änderungen für die einzelnen vorgegebenen Zeitpunkte durch einfache Interpolation bestimmen.
